\section{Definitions and Methods}
\label{sec:methods}

\subsection{The objects, in plain terms}

Everything in this paper is built from a finite-state Markov chain and a choice of what to
observe. Five objects carry all the results. We define each in ordinary terms, with the Kovacs
experiment as the running example, and give its source in the Six Birds Theory (SBT) corpus,
which we cite rather than re-derive \sbtcite{sbt-primer,sbt-foundations-i,sbt-foundations-ii,%
sbt-foundations-iii,sbt-foundations-iv,sbt-paper-a,sbt-paper-b,sbt-paper-c}.

\begin{description}[leftmargin=1.5em, itemsep=4pt]
  \item[Lens.] What is observed: a map $f$ from the microstate space $Z$ to a finite set of
    readout values. Our main lens is the energy readout $L_{\mathrm{energy}}$, the number of up
    spins; the ``thermodynamic lens'' $\Pi_{\mathrm{std}}$ is this readout together with its
    declared panel. In the laboratory the lens is the dilatometer. (SBT: a finite theory package,
    F-I \sbtcite{sbt-foundations-i}, Def.~1.)
  \item[Closure and the closure deficit.] A lens is \emph{closed} at horizon $\tau$, for a given
    law and a given distribution over microstates, if the readout now determines the
    distribution of the readout $\tau$ steps later, with nothing further gained from the
    microstate. The \emph{closure deficit}
    \[
      \mathrm{CD}_{\tau}(\Pi) \;=\; I\!\left(X_t;\, Y_{t+\tau} \,\middle|\, Y_t\right)
    \]
    is the conditional mutual information, in nats \cite{cover-thomas-2006}, between the
    microstate $X_t$ and the future readout $Y_{t+\tau}$ given the present readout $Y_t$. It is
    zero exactly when the $\tau$-step readout transition probabilities are the same from every
    \emph{positive-mass} microstate with the same readout. When the distribution has
    full support this is $\tau$-lumpability of the (kernel, lens) pair
    \cite{burke-rosenblatt-1958,kemeny-snell-1960,buchholz-1994}. With partial
    support, a zero deficit tests conditional independence only on the populated microstates at
    the declared time and horizon; zero-mass regions are never audited. A point-mass start on East
    $N=2$ has $\mathrm{CD}=0$ although the kernel is not lumpable, a regression-tested control.
    This does not by itself establish weak lumpability of the aggregated process, which is a
    statement about the whole readout process under an initial law
    \cite{rubino-sericola-1989}. Lumpability at one
    horizon does not imply it at others, so the ladder of horizons is declared. A positive
    deficit is an exact lower bound on what any autonomous law at that horizon, written in the
    lens's vocabulary, must discard. A readout that is not closed obeys a law with a memory
    kernel, in the projection-operator sense of \cite{zwanzig-1961,mori-1965}. (SBT: F-III
    \sbtcite{sbt-foundations-iii} T8; \sbtcite{sbt-paper-b} Def.~3.1.)
  \item[Packaging map and idempotence defect.] A second, coarser diagnostic. Evolve the chain
    $\tau$ steps, push the distribution through the lens, and lift it back to a microstate
    distribution by a declared rule (the ``prototype'' for each readout value):
    $E_{\tau,f}(\mu) = U_f(Q_f(\mu P^{\tau}))$. The \emph{idempotence defect} $\delta_{\tau,f}$
    is the largest total-variation distance, over starting microstates, between applying this map
    twice and applying it once. It is a saturation diagnostic and nothing more. A lumpable lens
    can still have $\delta_{\tau,f} > 0$: the unconstrained control at $N=2$, $\epsilon=1/2$,
    $\tau=1$ has exact defect $2/9$, a regression test. The defect is also distinct from the
    \emph{prototype-retention error}, the largest probability, over readout values, that the
    chain started from a value's prototype shows a different readout $\tau$ steps later; the
    retention error bounds the defect. GT1 reports the defect only as a ratio
    between two kernels, never as a closure criterion. (SBT: F-I Def.~9.)
  \item[Preparation histories, quotients, and witnesses.] A \emph{preparation history} is a
    protocol word (``equilibrate at $\epsilon_1$, quench to $\epsilon_2$, age $t_w$, jump to
    $\epsilon_3$'') together with the state distribution it produces. Over a declared finite
    catalog of histories and continuation protocols, the \emph{current quotient} $Q$ identifies
    histories with equal present lens readouts. The \emph{predictive quotient} $M$ identifies
    histories with equal future readout statistics under every declared continuation: the
    causal-state construction of computational mechanics
    \cite{crutchfield-young-1989,shalizi-crutchfield-2001,shalizi-moore-2025}, restricted to a
    declared finite catalog. When future equivalence implies current equivalence, $M$ refines $Q$
    and there is a comparison map $\pi: M \to Q$. This refinement is a premise, not automatic:
    the Python pipeline verifies it on every declared carrier before forming $\pi$, and in the
    Lean cores it follows because the identity continuation and the present events are among the
    declared futures. A \emph{witness}, or split pair, is two histories $h, h'$ that $Q$ identifies
    and $M$ separates: same now, different later (Fig.~\ref{fig:witness-concept}). Witnesses exist
    exactly when $\pi$ merges distinct $M$-classes. The Kovacs sample at its crossing moment and
    an equilibrated sample at the same volume are the laboratory picture of a witness pair.
    (SBT: \sbtcite{sbt-paper-a}, Defs.~2.2--2.4 and Thm.~3.5.)
  \item[Memory coordinates and buyback.] A witness pair can be resolved by adding coordinates to
    the lens. The \emph{buyback curve} records, for each budget of added scalars, which declared
    candidate coordinates resolve the pair. Each addition is a recorded, selected, non-unique
    choice (a ``lift'' in SBT terms). Fictive temperature is one such candidate. (SBT: F-II
    \texttt{prop:forgetting-lifts}.)
\end{description}

\begin{figure}[htbp]
  \centering
  \begin{tikzpicture}[
      hist/.style={draw, rounded corners=2pt, minimum width=22mm, minimum height=7mm,
        font=\small, fill=black!4},
      cls/.style={draw, circle, minimum size=8mm, font=\small, thick},
      grp/.style={draw, rounded corners=4pt, dashed, black!45},
      lab/.style={font=\footnotesize\itshape, text=black!65},
      ar/.style={-{Stealth[length=2mm]}, semithick}]
    \node[hist] (h1) at (0, 0.6) {history $h$};
    \node[hist] (h2) at (0,-0.6) {history $h'$};
    \node[cls, draw=orange!75!black] (m1) at (-4.2, 0.6) {$m$};
    \node[cls, draw=blue!65!black] (m2) at (-4.2,-0.6) {$m'$};
    \node[cls] (q) at (4.2, 0) {$q$};
    \node[grp, fit=(m1)(m2), inner sep=2.5mm] (M) {};
    \node[grp, fit=(q), inner sep=2.5mm, minimum height=24mm] (Q) {};
    \node[above=0.5mm of M, font=\small] {$M$: predictive};
    \node[above=0.5mm of Q, font=\small] {$Q$: current};
    \draw[ar] (h1) -- node[above, lab] {futures differ} (m1);
    \draw[ar] (h2) -- (m2);
    \draw[ar] (h1) -- node[above, lab, sloped] {readout now} (q);
    \draw[ar] (h2) -- (q);
    \draw[ar, black!70] (M.south) to[out=-25, in=-155] node[below, lab] {$\pi: M \to Q$ merges
      $m$ and $m'$} (Q.south);
  \end{tikzpicture}
  \caption{A witness pair. Two preparation histories have the same current readout (the
    current quotient $Q$ sends both to one class $q$) but different futures (the predictive
    quotient $M$ keeps them apart). The comparison map $\pi$ then merges two predictive classes,
    and no function of the current readout can predict the future. GT2 exhibits such a pair for the
    mean energy; GT3 and GT6 for the complete energy distribution, on a constructed pair and on
    an aging run respectively.}
  \label{fig:witness-concept}
\end{figure}

Four more terms recur. A \emph{lens-definable predicate} is a yes/no function of the readout
value; a static order parameter built from the lens is one of these. A \emph{shell} is a frozen
set of eligibility conditions under which a certificate is claimed: for GT6, the model and $N$,
the temperature schedule, the observation window, and a non-equilibration criterion, each checked
by a membership predicate with an exit counter that must read zero. A \emph{strict extension}
(SBT: a strict promotion) is the statement that the predictive description $M$ carries
information the declared present description $Q$ cannot. At its core is a
\emph{nonfactorization} statement: $M$ is not a function of $Q$. The framework's fuller filing
adds a checklist (a ``gate table'') and ablations (a ``knockout panel''). Finally, the \emph{six
roles} P1--P6 are the framework's fixed list of ways a description can fail to close. This paper
files only P4 (``the future depends on more than the present readout'') and P5 (``the packaging
of the present readout is not a closed projection''), and glosses them where they appear.

\paragraph{Claim grades and nonclaims.} Every result carries exactly one grade from the
framework's fidelity scale (F-II \sbtcite{sbt-foundations-ii}, \S9.3).
\emph{Theorem-grade (audited class)}: proved or audited unconditionally within the covered finite
scope, at the artifact's declared evidence type. \emph{Certificate-grade (shell-local)}, the
standard of paper [C] \sbtcite{sbt-paper-c}: valid on the frozen shell, with no broader-class
transfer. \emph{Run-only readout}: a measured number with no model fitted to it.
\emph{Recognition-grade}: a demonstration rather than a theorem. \emph{Control}: a check that a
method returns the known answer. Every artifact carries a machine-readable envelope
(\texttt{claim\_id}, \texttt{claim\_grade}, \texttt{quantifier\_domain}, \texttt{nonclaims})
validated against a schema in the reproducibility package. Nonclaims are not boilerplate; several
are load-bearing, and the full ledger is printed in Appendix~\ref{app:nonclaims}.

\subsection{Exact arithmetic and the freeze discipline}

Exact rational arithmetic (Python \texttt{fractions.Fraction}) carries the kernels, stationary
measures, witness gaps, the GT3 predicate sweep, the GT6 full-energy-law witness, the
trap-invariance lemma, the GT1 idempotence-defect ratio ($289/121$), and the lumpable zero control
(whose deficit is exactly $0$ because exact lumpability is verified symbolically). Quantities that
involve logarithms are certified rather than computed in floating point. GT1's deficits are
evaluated from exact joint and conditional probabilities, and every logarithm is enclosed between
outward-rounded rational bounds; the artifact declares \texttt{evidence\_type:
exact\_rational}, and every inequality of its margin is a statement about those
rational logarithm enclosures. The GT6 pressure bounds use the same device for exponentials and
logarithms. Each artifact declares its evidence type, and a grade rests on the audited-class
quantification plus that declared evidence type. Floating point (\texttt{float64\_deterministic})
appears only in explicitly labeled diagnostics: printed decimal values, spectral gaps, windowed
deficit diagnostics, GT7 Monte Carlo spot checks, the finite pressure ladder and the GB1
three-state demonstration. Controls accompany every certificate, including the constructed
negative controls of the triage panel (\S\ref{sec:gt3-gt4}).

The paper's evidence base is frozen by sha256 content hashes over canonical JSON. Each freeze
artifact carries a \texttt{content\_hash}, and every cross-artifact citation pins the target's
hash, so regenerating any artifact invalidates, detectably, every citation of it. A
repository-wide test suite enforces these bindings. The paper's own claim citations, the bracketed
row labels such as [GT1], are pinned the same way: each carries a hidden artifact hash that is
checked against the claims ledger and the current artifact bytes at build time. Content hashes
establish \emph{which} bytes were frozen; they do not, by themselves, establish \emph{when}
(\S\ref{sec:gt7}).

\subsection{What the reader can check}

The claims ledger (\texttt{paper/claims\_ledger.json}) is the paper's skeleton: one row per claim
with its final wording, grade, quantifier domain, anchors into the SBT corpus, artifact hashes and
nonclaims. Each results section transcribes its ledger rows into prose; a discrepancy between
prose and ledger is a bug in the prose. A reader with the repository can re-run every certificate
from a clean checkout: the test suite, the Lean build with its per-theorem axiom audit, and the
three checkers (claims ledger, GT6 filing, empirical bridges).
